\documentclass{article}
\usepackage{amsmath}  % برای استفاده از توابع ریاضی
\usepackage{amssymb}
\usepackage{longtable}
\usepackage{amsthm}

\usepackage{array}
\usepackage{booktabs}
\usepackage{multirow}
\usepackage{multicol}
\usepackage{graphicx}
\usepackage{subcaption}
\usepackage{hyperref}
\usepackage{siunitx}
\sisetup{group-separator = {,}}
\usepackage{xepersian}
\settextfont[Scale=1]{XB Niloofar}
\newtheorem{definition}{تعریف}[section]

\begin{document}
%	\documentclass[a4paper 11px twocolumn]{article}
% تنظیمات و پکیج های مورد نیاز در این فایل قرار دارد. برای تنظیم و افزودن پکیج به این فایل مراجعه کنید.
\usepackage{amsmath, amssymb, amsthm} % بسته‌های ریاضی و قضایای امسال
\usepackage[all]{xy} % برای ترسیم نمودارها
\usepackage[left=1.3cm, right=1.3cm, top=1cm, bottom=1.3cm, nohead, hmargin=2cm,vmargin=2.5cm]{geometry} % تنظیمات حاشیه‌ها
\usepackage{graphicx} % برای درج تصاویر
\usepackage{fancyhdr} % برای تنظیم سربرگ و پاورقی صفحات
\usepackage{fontspec} % برای انتخاب فونتها
\usepackage{framed} % برای ایجاد قاب برای متن
\usepackage{ } % برای شمارش تعداد صفحات
\usepackage{cite} % برای ارجاع به منابع
\usepackage{pdfpages} % برای درج صفحات PDF
\usepackage{ifthen} % برای شرطی کردن تنظیمات
\usepackage{xcolor, colortbl} % برای استفاده از رنگها
\usepackage{pifont} % برای استفاده از نمادها

\usepackage{enumitem} % برای سفارشی‌سازی محیط‌های شماره‌گذاری
\usepackage{stackengine} % برای ترتیب متون و اشیاء
\usepackage{tabularx} % برای جداول با عرض متغیر
\usepackage{multirow} % برای ترکیب سلول‌های چند ردیفی در جدول
\usepackage{fourier} % برای فونتهای تایمز
\usepackage{array} % برای سفارشی‌سازی جداول
\usepackage{makecell} % برای سفارشی‌سازی سلول‌ها در جدول
\usepackage{ptext} % تولید متن تصادفی
\usepackage{hyperref} % برای ایجاد لینکهای رنگی
\usepackage[nottoc]{tocbibind} % برای اضافه کردن فهرست مطالب به فهرست مطالب
\usepackage{bidipoem} % برای شعر فارسی
\usepackage{tikz} % برای ترسیم نمودارهای تعاملی
\usepackage{atbegshi} % برای کنترل تصاویر و نمودارها
\usetikzlibrary{calc} % کتابخانه‌ای برای ترسیم نمودارهای TikZ


\usepackage{fancyhdr} % Custom headers and footers


% چنانچه قصد پرینت گرفتن نوشته خود را دارید، خط بالا را غیرفعال و  از دستور زیر استفاده کنید چون در صورت استفاده از دستور زیر‌‌، 
% لینک‌ها به رنگ سیاه ظاهر خواهند شد که برای پرینت گرفتن، مناسب‌تر است
%\usepackage[pagebackref=false]{hyperref}


% بسته‌ لازم برای تنظیم سربرگ‌ها
% بسته‌ای برای ظاهر شدن «مراجع» و «نمایه» در فهرست مطالب

% دستورات مربوط به ایجاد نمایه
%\usepackage{makeidx}
%\makeindex
%%%%%%%%%%%%%%%%%%%%%%%%%%
% فراخوانی بسته زی‌پرشین و تعریف قلم فارسی و انگلیسی
\usepackage{xepersian}
\settextfont[Scale=1]{XB Niloofar}
%%%%%%%%%%%%%%%%%%%%%%%%%%
% چنانچه می‌خواهید اعداد در فرمول‌ها، انگلیسی باشد، خط زیر را غیرفعال کنید
\setdigitfont[Scale=1.2]{XB Niloofar}

\defpersianfont\nas[Scale=2]{IranNastaliq}

\defpersianfont\titr[Scale=1.5]{B Titr Bold}

%%%%%%%%%%%%%%%%%%%%%%%%%%


%%%%%%%%%%%%%%%%%%%%%%%%%%
% تعریف و نحوه ظاهر شدن عنوان قضیه‌ها، تعریف‌ها، مثال‌ها و ...
\theoremstyle{definition}
\newtheorem{definition}{تعریف}[section]
\theoremstyle{theorem}
%\newtheorem{s}[definition]{}
\newtheorem{theorem}[definition]{قضیه}
\newtheorem{lemma}[definition]{لم}
\newtheorem{proposition}[definition]{گزاره}
\newtheorem{corollary}[definition]{نتیجه}
\newtheorem{remark}[definition]{تبصره}
\theoremstyle{definition}
\newtheorem{example}[definition]{مثال}
\newtheorem{notation}[definition]{نمادگذاری}
\newtheorem{definitiont}[definition]{تعریف و نمادگذاری}
\newtheorem{obs}[definition]{ملاحظه}
%%%%%%%%%%%%%%%%%%%%%%%%%%%%

\begin{document}
% مشخصات گزارش را در این قسمت تغییر بدهید.
\begin{titlepage}
	\setlength{\arrayrulewidth}{0.4mm}
	\newcommand{\HRule}{\rule{\linewidth}{0.5mm}}
	\begin{center}
		% logo
		\includegraphics[width=0.5\textwidth]{./Fig/logo.png}~\\[4cm]
		
		\HRule \\[0.4cm]
		{
			\Huge\textbf{عنوان گزارش} \\
			[0.4cm]
			
			\large \bfseries
			نام درس \\ 
			[0.4cm]}
		\HRule \\[4cm]
		
		\large\textbf{نویسنده: }\\[0.2cm]
		نام و نام خوانوادگی 
		\\[3cm]
		
		\vfill
		
		{\large \today}
		
	\end{center}
\end{titlepage}

% متن مربوط به گزارش را در این قسمت قرار دهید.
\twocolumn
\ptext[0-10]


% در صورت نیاز به مراجع و پیوست از دو خط و فایل های مربوطه استفاده فرمایید.
% پیوست خود را در این قسمت وارد کنید.
\appendix
\onecolumn
% شروع محیط پیوست 
\section{پیوست}
\ptext[10-20]

% مراجع خود را در این قسمت وارد کنید.
% دستوری برای کوچک کردن اندازه فونت‌ها 
\small
% شروع محیط مراجع 

\renewcommand\bibname{مراجع}
\setLTRbibitems
% برای ارجاع دهی خودکار محتوای فایل bibtex مراجع را در فایل Refrence/myreferences.bib قرار داده و ذخیره کنید.  و یا میتوانید دو خط زیر را کامنت کرده و از دستورات بعدی ان بزرای ارجاع دهی استفاده کنید.
\bibliographystyle{IEEEtran}
\bibliography{Refrence/myreferences.bib}
%\begin{thebibliography}{99}
%
%\resetlatinfont
%% تذکر: چنانچه از نسخه‌های قدیمی زی‌پرشین و بی‌دی استفاده می‌کنید، هنگام استفاده از سه دستور بالا، با خطا 
%% مواجه می‌شوید. بنابراین باید زی‌پرشین و بی‌دی خود را آپدیت کنید.
%% چنانچه مرجع فارسی هم دارید باید یا از بسته Persian-bib استفاده کنید و یا راهنمای bidi را مطالعه کنید.
%
%\bibitem {TE} { \sc S. T. Aldrich, E.  Enochs, J. R. Garc\'{i}a Rozas, L. Oyonarte,} {\sl Covers and envelopes in Grothendieck categories: flat covers of complexes with applications,} J. Algebra {\bf 243} (2001), 615-630.
%
%\bibitem {AJL}{ \sc L. Alonso, A. Jerem\'{i}oas, J. Lipman,} {\sl  Duality and flat base change on formal schemes,}
%in “Studies in Duality on Noetherian Formal Schemes and Non-Noetherian Ordinary
%Schemes,” Contemp. Math. {\bf 244}, Amer. Math. Soc., 1999, 3-90.
%
%\bibitem {AF} {\sc F. W. Anderson, K. R. Fuller,} {\sl Rings and categories of modules,} Graduate Texts in Mathematics 13, Springer-Verlag, Berlin-Heidelberg-New York, 1974.
%
%\bibitem {AEHS} {\sc J. Asadollahi, H. Eshraghi, R. Hafezi, Sh. Salarian,} {\sl On the homotopy categories of projective and injective representations of quivers,} J. Algebra {\bf 346} (2011), 101-115.
%
%
%
%
%\end{thebibliography}
%
%

\end{document}
%	\documentclass[a4paper 11px twocolumn]{article}
% تنظیمات و پکیج های مورد نیاز در این فایل قرار دارد. برای تنظیم و افزودن پکیج به این فایل مراجعه کنید.
\usepackage{amsmath, amssymb, amsthm} % بسته‌های ریاضی و قضایای امسال
\usepackage[all]{xy} % برای ترسیم نمودارها
\usepackage[left=1.3cm, right=1.3cm, top=1cm, bottom=1.3cm, nohead, hmargin=2cm,vmargin=2.5cm]{geometry} % تنظیمات حاشیه‌ها
\usepackage{graphicx} % برای درج تصاویر
\usepackage{fancyhdr} % برای تنظیم سربرگ و پاورقی صفحات
\usepackage{fontspec} % برای انتخاب فونتها
\usepackage{framed} % برای ایجاد قاب برای متن
\usepackage{ } % برای شمارش تعداد صفحات
\usepackage{cite} % برای ارجاع به منابع
\usepackage{pdfpages} % برای درج صفحات PDF
\usepackage{ifthen} % برای شرطی کردن تنظیمات
\usepackage{xcolor, colortbl} % برای استفاده از رنگها
\usepackage{pifont} % برای استفاده از نمادها

\usepackage{enumitem} % برای سفارشی‌سازی محیط‌های شماره‌گذاری
\usepackage{stackengine} % برای ترتیب متون و اشیاء
\usepackage{tabularx} % برای جداول با عرض متغیر
\usepackage{multirow} % برای ترکیب سلول‌های چند ردیفی در جدول
\usepackage{fourier} % برای فونتهای تایمز
\usepackage{array} % برای سفارشی‌سازی جداول
\usepackage{makecell} % برای سفارشی‌سازی سلول‌ها در جدول
\usepackage{ptext} % تولید متن تصادفی
\usepackage{hyperref} % برای ایجاد لینکهای رنگی
\usepackage[nottoc]{tocbibind} % برای اضافه کردن فهرست مطالب به فهرست مطالب
\usepackage{bidipoem} % برای شعر فارسی
\usepackage{tikz} % برای ترسیم نمودارهای تعاملی
\usepackage{atbegshi} % برای کنترل تصاویر و نمودارها
\usetikzlibrary{calc} % کتابخانه‌ای برای ترسیم نمودارهای TikZ


\usepackage{fancyhdr} % Custom headers and footers


% چنانچه قصد پرینت گرفتن نوشته خود را دارید، خط بالا را غیرفعال و  از دستور زیر استفاده کنید چون در صورت استفاده از دستور زیر‌‌، 
% لینک‌ها به رنگ سیاه ظاهر خواهند شد که برای پرینت گرفتن، مناسب‌تر است
%\usepackage[pagebackref=false]{hyperref}


% بسته‌ لازم برای تنظیم سربرگ‌ها
% بسته‌ای برای ظاهر شدن «مراجع» و «نمایه» در فهرست مطالب

% دستورات مربوط به ایجاد نمایه
%\usepackage{makeidx}
%\makeindex
%%%%%%%%%%%%%%%%%%%%%%%%%%
% فراخوانی بسته زی‌پرشین و تعریف قلم فارسی و انگلیسی
\usepackage{xepersian}
\settextfont[Scale=1]{XB Niloofar}
%%%%%%%%%%%%%%%%%%%%%%%%%%
% چنانچه می‌خواهید اعداد در فرمول‌ها، انگلیسی باشد، خط زیر را غیرفعال کنید
\setdigitfont[Scale=1.2]{XB Niloofar}

\defpersianfont\nas[Scale=2]{IranNastaliq}

\defpersianfont\titr[Scale=1.5]{B Titr Bold}

%%%%%%%%%%%%%%%%%%%%%%%%%%


%%%%%%%%%%%%%%%%%%%%%%%%%%
% تعریف و نحوه ظاهر شدن عنوان قضیه‌ها، تعریف‌ها، مثال‌ها و ...
\theoremstyle{definition}
\newtheorem{definition}{تعریف}[section]
\theoremstyle{theorem}
%\newtheorem{s}[definition]{}
\newtheorem{theorem}[definition]{قضیه}
\newtheorem{lemma}[definition]{لم}
\newtheorem{proposition}[definition]{گزاره}
\newtheorem{corollary}[definition]{نتیجه}
\newtheorem{remark}[definition]{تبصره}
\theoremstyle{definition}
\newtheorem{example}[definition]{مثال}
\newtheorem{notation}[definition]{نمادگذاری}
\newtheorem{definitiont}[definition]{تعریف و نمادگذاری}
\newtheorem{obs}[definition]{ملاحظه}
%%%%%%%%%%%%%%%%%%%%%%%%%%%%

\begin{document}
% مشخصات گزارش را در این قسمت تغییر بدهید.
\begin{titlepage}
	\setlength{\arrayrulewidth}{0.4mm}
	\newcommand{\HRule}{\rule{\linewidth}{0.5mm}}
	\begin{center}
		% logo
		\includegraphics[width=0.5\textwidth]{./Fig/logo.png}~\\[4cm]
		
		\HRule \\[0.4cm]
		{
			\Huge\textbf{عنوان گزارش} \\
			[0.4cm]
			
			\large \bfseries
			نام درس \\ 
			[0.4cm]}
		\HRule \\[4cm]
		
		\large\textbf{نویسنده: }\\[0.2cm]
		نام و نام خوانوادگی 
		\\[3cm]
		
		\vfill
		
		{\large \today}
		
	\end{center}
\end{titlepage}

% متن مربوط به گزارش را در این قسمت قرار دهید.
\twocolumn
\ptext[0-10]


% در صورت نیاز به مراجع و پیوست از دو خط و فایل های مربوطه استفاده فرمایید.
% پیوست خود را در این قسمت وارد کنید.
\appendix
\onecolumn
% شروع محیط پیوست 
\section{پیوست}
\ptext[10-20]

% مراجع خود را در این قسمت وارد کنید.
% دستوری برای کوچک کردن اندازه فونت‌ها 
\small
% شروع محیط مراجع 

\renewcommand\bibname{مراجع}
\setLTRbibitems
% برای ارجاع دهی خودکار محتوای فایل bibtex مراجع را در فایل Refrence/myreferences.bib قرار داده و ذخیره کنید.  و یا میتوانید دو خط زیر را کامنت کرده و از دستورات بعدی ان بزرای ارجاع دهی استفاده کنید.
\bibliographystyle{IEEEtran}
\bibliography{Refrence/myreferences.bib}
%\begin{thebibliography}{99}
%
%\resetlatinfont
%% تذکر: چنانچه از نسخه‌های قدیمی زی‌پرشین و بی‌دی استفاده می‌کنید، هنگام استفاده از سه دستور بالا، با خطا 
%% مواجه می‌شوید. بنابراین باید زی‌پرشین و بی‌دی خود را آپدیت کنید.
%% چنانچه مرجع فارسی هم دارید باید یا از بسته Persian-bib استفاده کنید و یا راهنمای bidi را مطالعه کنید.
%
%\bibitem {TE} { \sc S. T. Aldrich, E.  Enochs, J. R. Garc\'{i}a Rozas, L. Oyonarte,} {\sl Covers and envelopes in Grothendieck categories: flat covers of complexes with applications,} J. Algebra {\bf 243} (2001), 615-630.
%
%\bibitem {AJL}{ \sc L. Alonso, A. Jerem\'{i}oas, J. Lipman,} {\sl  Duality and flat base change on formal schemes,}
%in “Studies in Duality on Noetherian Formal Schemes and Non-Noetherian Ordinary
%Schemes,” Contemp. Math. {\bf 244}, Amer. Math. Soc., 1999, 3-90.
%
%\bibitem {AF} {\sc F. W. Anderson, K. R. Fuller,} {\sl Rings and categories of modules,} Graduate Texts in Mathematics 13, Springer-Verlag, Berlin-Heidelberg-New York, 1974.
%
%\bibitem {AEHS} {\sc J. Asadollahi, H. Eshraghi, R. Hafezi, Sh. Salarian,} {\sl On the homotopy categories of projective and injective representations of quivers,} J. Algebra {\bf 346} (2011), 101-115.
%
%
%
%
%\end{thebibliography}
%
%

\end{document}
%	\documentclass[a4paper 11px twocolumn]{article}
% تنظیمات و پکیج های مورد نیاز در این فایل قرار دارد. برای تنظیم و افزودن پکیج به این فایل مراجعه کنید.
\usepackage{amsmath, amssymb, amsthm} % بسته‌های ریاضی و قضایای امسال
\usepackage[all]{xy} % برای ترسیم نمودارها
\usepackage[left=1.3cm, right=1.3cm, top=1cm, bottom=1.3cm, nohead, hmargin=2cm,vmargin=2.5cm]{geometry} % تنظیمات حاشیه‌ها
\usepackage{graphicx} % برای درج تصاویر
\usepackage{fancyhdr} % برای تنظیم سربرگ و پاورقی صفحات
\usepackage{fontspec} % برای انتخاب فونتها
\usepackage{framed} % برای ایجاد قاب برای متن
\usepackage{ } % برای شمارش تعداد صفحات
\usepackage{cite} % برای ارجاع به منابع
\usepackage{pdfpages} % برای درج صفحات PDF
\usepackage{ifthen} % برای شرطی کردن تنظیمات
\usepackage{xcolor, colortbl} % برای استفاده از رنگها
\usepackage{pifont} % برای استفاده از نمادها

\usepackage{enumitem} % برای سفارشی‌سازی محیط‌های شماره‌گذاری
\usepackage{stackengine} % برای ترتیب متون و اشیاء
\usepackage{tabularx} % برای جداول با عرض متغیر
\usepackage{multirow} % برای ترکیب سلول‌های چند ردیفی در جدول
\usepackage{fourier} % برای فونتهای تایمز
\usepackage{array} % برای سفارشی‌سازی جداول
\usepackage{makecell} % برای سفارشی‌سازی سلول‌ها در جدول
\usepackage{ptext} % تولید متن تصادفی
\usepackage{hyperref} % برای ایجاد لینکهای رنگی
\usepackage[nottoc]{tocbibind} % برای اضافه کردن فهرست مطالب به فهرست مطالب
\usepackage{bidipoem} % برای شعر فارسی
\usepackage{tikz} % برای ترسیم نمودارهای تعاملی
\usepackage{atbegshi} % برای کنترل تصاویر و نمودارها
\usetikzlibrary{calc} % کتابخانه‌ای برای ترسیم نمودارهای TikZ


\usepackage{fancyhdr} % Custom headers and footers


% چنانچه قصد پرینت گرفتن نوشته خود را دارید، خط بالا را غیرفعال و  از دستور زیر استفاده کنید چون در صورت استفاده از دستور زیر‌‌، 
% لینک‌ها به رنگ سیاه ظاهر خواهند شد که برای پرینت گرفتن، مناسب‌تر است
%\usepackage[pagebackref=false]{hyperref}


% بسته‌ لازم برای تنظیم سربرگ‌ها
% بسته‌ای برای ظاهر شدن «مراجع» و «نمایه» در فهرست مطالب

% دستورات مربوط به ایجاد نمایه
%\usepackage{makeidx}
%\makeindex
%%%%%%%%%%%%%%%%%%%%%%%%%%
% فراخوانی بسته زی‌پرشین و تعریف قلم فارسی و انگلیسی
\usepackage{xepersian}
\settextfont[Scale=1]{XB Niloofar}
%%%%%%%%%%%%%%%%%%%%%%%%%%
% چنانچه می‌خواهید اعداد در فرمول‌ها، انگلیسی باشد، خط زیر را غیرفعال کنید
\setdigitfont[Scale=1.2]{XB Niloofar}

\defpersianfont\nas[Scale=2]{IranNastaliq}

\defpersianfont\titr[Scale=1.5]{B Titr Bold}

%%%%%%%%%%%%%%%%%%%%%%%%%%


%%%%%%%%%%%%%%%%%%%%%%%%%%
% تعریف و نحوه ظاهر شدن عنوان قضیه‌ها، تعریف‌ها، مثال‌ها و ...
\theoremstyle{definition}
\newtheorem{definition}{تعریف}[section]
\theoremstyle{theorem}
%\newtheorem{s}[definition]{}
\newtheorem{theorem}[definition]{قضیه}
\newtheorem{lemma}[definition]{لم}
\newtheorem{proposition}[definition]{گزاره}
\newtheorem{corollary}[definition]{نتیجه}
\newtheorem{remark}[definition]{تبصره}
\theoremstyle{definition}
\newtheorem{example}[definition]{مثال}
\newtheorem{notation}[definition]{نمادگذاری}
\newtheorem{definitiont}[definition]{تعریف و نمادگذاری}
\newtheorem{obs}[definition]{ملاحظه}
%%%%%%%%%%%%%%%%%%%%%%%%%%%%

\begin{document}
% مشخصات گزارش را در این قسمت تغییر بدهید.
\begin{titlepage}
	\setlength{\arrayrulewidth}{0.4mm}
	\newcommand{\HRule}{\rule{\linewidth}{0.5mm}}
	\begin{center}
		% logo
		\includegraphics[width=0.5\textwidth]{./Fig/logo.png}~\\[4cm]
		
		\HRule \\[0.4cm]
		{
			\Huge\textbf{عنوان گزارش} \\
			[0.4cm]
			
			\large \bfseries
			نام درس \\ 
			[0.4cm]}
		\HRule \\[4cm]
		
		\large\textbf{نویسنده: }\\[0.2cm]
		نام و نام خوانوادگی 
		\\[3cm]
		
		\vfill
		
		{\large \today}
		
	\end{center}
\end{titlepage}

% متن مربوط به گزارش را در این قسمت قرار دهید.
\twocolumn
\ptext[0-10]


% در صورت نیاز به مراجع و پیوست از دو خط و فایل های مربوطه استفاده فرمایید.
% پیوست خود را در این قسمت وارد کنید.
\appendix
\onecolumn
% شروع محیط پیوست 
\section{پیوست}
\ptext[10-20]

% مراجع خود را در این قسمت وارد کنید.
% دستوری برای کوچک کردن اندازه فونت‌ها 
\small
% شروع محیط مراجع 

\renewcommand\bibname{مراجع}
\setLTRbibitems
% برای ارجاع دهی خودکار محتوای فایل bibtex مراجع را در فایل Refrence/myreferences.bib قرار داده و ذخیره کنید.  و یا میتوانید دو خط زیر را کامنت کرده و از دستورات بعدی ان بزرای ارجاع دهی استفاده کنید.
\bibliographystyle{IEEEtran}
\bibliography{Refrence/myreferences.bib}
%\begin{thebibliography}{99}
%
%\resetlatinfont
%% تذکر: چنانچه از نسخه‌های قدیمی زی‌پرشین و بی‌دی استفاده می‌کنید، هنگام استفاده از سه دستور بالا، با خطا 
%% مواجه می‌شوید. بنابراین باید زی‌پرشین و بی‌دی خود را آپدیت کنید.
%% چنانچه مرجع فارسی هم دارید باید یا از بسته Persian-bib استفاده کنید و یا راهنمای bidi را مطالعه کنید.
%
%\bibitem {TE} { \sc S. T. Aldrich, E.  Enochs, J. R. Garc\'{i}a Rozas, L. Oyonarte,} {\sl Covers and envelopes in Grothendieck categories: flat covers of complexes with applications,} J. Algebra {\bf 243} (2001), 615-630.
%
%\bibitem {AJL}{ \sc L. Alonso, A. Jerem\'{i}oas, J. Lipman,} {\sl  Duality and flat base change on formal schemes,}
%in “Studies in Duality on Noetherian Formal Schemes and Non-Noetherian Ordinary
%Schemes,” Contemp. Math. {\bf 244}, Amer. Math. Soc., 1999, 3-90.
%
%\bibitem {AF} {\sc F. W. Anderson, K. R. Fuller,} {\sl Rings and categories of modules,} Graduate Texts in Mathematics 13, Springer-Verlag, Berlin-Heidelberg-New York, 1974.
%
%\bibitem {AEHS} {\sc J. Asadollahi, H. Eshraghi, R. Hafezi, Sh. Salarian,} {\sl On the homotopy categories of projective and injective representations of quivers,} J. Algebra {\bf 346} (2011), 101-115.
%
%
%
%
%\end{thebibliography}
%
%

\end{document}
%	\documentclass[a4paper 11px twocolumn]{article}
% تنظیمات و پکیج های مورد نیاز در این فایل قرار دارد. برای تنظیم و افزودن پکیج به این فایل مراجعه کنید.
\usepackage{amsmath, amssymb, amsthm} % بسته‌های ریاضی و قضایای امسال
\usepackage[all]{xy} % برای ترسیم نمودارها
\usepackage[left=1.3cm, right=1.3cm, top=1cm, bottom=1.3cm, nohead, hmargin=2cm,vmargin=2.5cm]{geometry} % تنظیمات حاشیه‌ها
\usepackage{graphicx} % برای درج تصاویر
\usepackage{fancyhdr} % برای تنظیم سربرگ و پاورقی صفحات
\usepackage{fontspec} % برای انتخاب فونتها
\usepackage{framed} % برای ایجاد قاب برای متن
\usepackage{ } % برای شمارش تعداد صفحات
\usepackage{cite} % برای ارجاع به منابع
\usepackage{pdfpages} % برای درج صفحات PDF
\usepackage{ifthen} % برای شرطی کردن تنظیمات
\usepackage{xcolor, colortbl} % برای استفاده از رنگها
\usepackage{pifont} % برای استفاده از نمادها

\usepackage{enumitem} % برای سفارشی‌سازی محیط‌های شماره‌گذاری
\usepackage{stackengine} % برای ترتیب متون و اشیاء
\usepackage{tabularx} % برای جداول با عرض متغیر
\usepackage{multirow} % برای ترکیب سلول‌های چند ردیفی در جدول
\usepackage{fourier} % برای فونتهای تایمز
\usepackage{array} % برای سفارشی‌سازی جداول
\usepackage{makecell} % برای سفارشی‌سازی سلول‌ها در جدول
\usepackage{ptext} % تولید متن تصادفی
\usepackage{hyperref} % برای ایجاد لینکهای رنگی
\usepackage[nottoc]{tocbibind} % برای اضافه کردن فهرست مطالب به فهرست مطالب
\usepackage{bidipoem} % برای شعر فارسی
\usepackage{tikz} % برای ترسیم نمودارهای تعاملی
\usepackage{atbegshi} % برای کنترل تصاویر و نمودارها
\usetikzlibrary{calc} % کتابخانه‌ای برای ترسیم نمودارهای TikZ


\usepackage{fancyhdr} % Custom headers and footers


% چنانچه قصد پرینت گرفتن نوشته خود را دارید، خط بالا را غیرفعال و  از دستور زیر استفاده کنید چون در صورت استفاده از دستور زیر‌‌، 
% لینک‌ها به رنگ سیاه ظاهر خواهند شد که برای پرینت گرفتن، مناسب‌تر است
%\usepackage[pagebackref=false]{hyperref}


% بسته‌ لازم برای تنظیم سربرگ‌ها
% بسته‌ای برای ظاهر شدن «مراجع» و «نمایه» در فهرست مطالب

% دستورات مربوط به ایجاد نمایه
%\usepackage{makeidx}
%\makeindex
%%%%%%%%%%%%%%%%%%%%%%%%%%
% فراخوانی بسته زی‌پرشین و تعریف قلم فارسی و انگلیسی
\usepackage{xepersian}
\settextfont[Scale=1]{XB Niloofar}
%%%%%%%%%%%%%%%%%%%%%%%%%%
% چنانچه می‌خواهید اعداد در فرمول‌ها، انگلیسی باشد، خط زیر را غیرفعال کنید
\setdigitfont[Scale=1.2]{XB Niloofar}

\defpersianfont\nas[Scale=2]{IranNastaliq}

\defpersianfont\titr[Scale=1.5]{B Titr Bold}

%%%%%%%%%%%%%%%%%%%%%%%%%%


%%%%%%%%%%%%%%%%%%%%%%%%%%
% تعریف و نحوه ظاهر شدن عنوان قضیه‌ها، تعریف‌ها، مثال‌ها و ...
\theoremstyle{definition}
\newtheorem{definition}{تعریف}[section]
\theoremstyle{theorem}
%\newtheorem{s}[definition]{}
\newtheorem{theorem}[definition]{قضیه}
\newtheorem{lemma}[definition]{لم}
\newtheorem{proposition}[definition]{گزاره}
\newtheorem{corollary}[definition]{نتیجه}
\newtheorem{remark}[definition]{تبصره}
\theoremstyle{definition}
\newtheorem{example}[definition]{مثال}
\newtheorem{notation}[definition]{نمادگذاری}
\newtheorem{definitiont}[definition]{تعریف و نمادگذاری}
\newtheorem{obs}[definition]{ملاحظه}
%%%%%%%%%%%%%%%%%%%%%%%%%%%%

\begin{document}
% مشخصات گزارش را در این قسمت تغییر بدهید.
\begin{titlepage}
	\setlength{\arrayrulewidth}{0.4mm}
	\newcommand{\HRule}{\rule{\linewidth}{0.5mm}}
	\begin{center}
		% logo
		\includegraphics[width=0.5\textwidth]{./Fig/logo.png}~\\[4cm]
		
		\HRule \\[0.4cm]
		{
			\Huge\textbf{عنوان گزارش} \\
			[0.4cm]
			
			\large \bfseries
			نام درس \\ 
			[0.4cm]}
		\HRule \\[4cm]
		
		\large\textbf{نویسنده: }\\[0.2cm]
		نام و نام خوانوادگی 
		\\[3cm]
		
		\vfill
		
		{\large \today}
		
	\end{center}
\end{titlepage}

% متن مربوط به گزارش را در این قسمت قرار دهید.
\twocolumn
\ptext[0-10]


% در صورت نیاز به مراجع و پیوست از دو خط و فایل های مربوطه استفاده فرمایید.
% پیوست خود را در این قسمت وارد کنید.
\appendix
\onecolumn
% شروع محیط پیوست 
\section{پیوست}
\ptext[10-20]

% مراجع خود را در این قسمت وارد کنید.
% دستوری برای کوچک کردن اندازه فونت‌ها 
\small
% شروع محیط مراجع 

\renewcommand\bibname{مراجع}
\setLTRbibitems
% برای ارجاع دهی خودکار محتوای فایل bibtex مراجع را در فایل Refrence/myreferences.bib قرار داده و ذخیره کنید.  و یا میتوانید دو خط زیر را کامنت کرده و از دستورات بعدی ان بزرای ارجاع دهی استفاده کنید.
\bibliographystyle{IEEEtran}
\bibliography{Refrence/myreferences.bib}
%\begin{thebibliography}{99}
%
%\resetlatinfont
%% تذکر: چنانچه از نسخه‌های قدیمی زی‌پرشین و بی‌دی استفاده می‌کنید، هنگام استفاده از سه دستور بالا، با خطا 
%% مواجه می‌شوید. بنابراین باید زی‌پرشین و بی‌دی خود را آپدیت کنید.
%% چنانچه مرجع فارسی هم دارید باید یا از بسته Persian-bib استفاده کنید و یا راهنمای bidi را مطالعه کنید.
%
%\bibitem {TE} { \sc S. T. Aldrich, E.  Enochs, J. R. Garc\'{i}a Rozas, L. Oyonarte,} {\sl Covers and envelopes in Grothendieck categories: flat covers of complexes with applications,} J. Algebra {\bf 243} (2001), 615-630.
%
%\bibitem {AJL}{ \sc L. Alonso, A. Jerem\'{i}oas, J. Lipman,} {\sl  Duality and flat base change on formal schemes,}
%in “Studies in Duality on Noetherian Formal Schemes and Non-Noetherian Ordinary
%Schemes,” Contemp. Math. {\bf 244}, Amer. Math. Soc., 1999, 3-90.
%
%\bibitem {AF} {\sc F. W. Anderson, K. R. Fuller,} {\sl Rings and categories of modules,} Graduate Texts in Mathematics 13, Springer-Verlag, Berlin-Heidelberg-New York, 1974.
%
%\bibitem {AEHS} {\sc J. Asadollahi, H. Eshraghi, R. Hafezi, Sh. Salarian,} {\sl On the homotopy categories of projective and injective representations of quivers,} J. Algebra {\bf 346} (2011), 101-115.
%
%
%
%
%\end{thebibliography}
%
%

\end{document}
%	\small
% شروع محیط مراجع 

\renewcommand\bibname{مراجع}
\setLTRbibitems

\bibliographystyle{IEEEtran}
\bibliography{./Refrence/myreferences.bib}
%\bibliography{./Refrence/otherreferences.bib}
		
\end{document}

